\documentclass[11pt]{article}
\usepackage{amsmath, amsthm, amssymb}
\usepackage{mathtools}
\usepackage[margin=1in]{geometry}
\usepackage{hyperref}

\theoremstyle{plain}
\newtheorem{theorem}{Theorem}
\newtheorem{lemma}[theorem]{Lemma}
\newtheorem{proposition}[theorem]{Proposition}
\newtheorem{corollary}[theorem]{Corollary}

\theoremstyle{definition}
\newtheorem{definition}[theorem]{Definition}
\newtheorem{example}[theorem]{Example}

\theoremstyle{remark}
\newtheorem{remark}[theorem]{Remark}

\newcommand{\Z}{\mathbb{Z}}
\newcommand{\N}{\mathbb{N}}
\newcommand{\Q}{\mathbb{Q}}
\newcommand{\R}{\mathbb{R}}
\newcommand{\Sym}{\mathfrak{S}}
\newcommand{\Hq}{H_q}
\newcommand{\Pibs}{\Pi_q}
\newcommand{\Vlam}{V_\lambda}
\newcommand{\trace}{\mathrm{tr}}

\title{Non-negativity of the trace core \\ for the staircase Bott--Samelson operator \\ on Hecke cell modules}
\author{Clio}
\date{24 April 2026}

\begin{document}
\maketitle

\begin{abstract}
Let $\Hq = \Hq(\Sym_n)$ be the Hecke algebra of the symmetric group with parameter $q$, and let $\Pibs = \prod_{j=1}^{\binom{n}{2}}(1 + T_{i_j})$ be the Bott--Samelson product along the staircase reduced word $\underline{w}_0 = s_1 \cdot s_2 s_1 \cdot s_3 s_2 s_1 \cdots s_{n-1} \cdots s_1$ for the longest element. For any partition $\lambda \vdash n$ let $V_\lambda$ denote the corresponding cell module. Then $\sigma_1(\Vlam) := \trace(\Pibs|_{\Vlam}) \in \Z[q]$ is palindromic (Apr 23 theorem) and non-negative (by Kazhdan--Lusztig positivity). We prove:
\begin{quote}
\itshape If $\sigma_1(\Vlam) = q^{a_\lambda}(1+q)^{b_\lambda} P_\lambda(q)$ with $(a_\lambda, b_\lambda)$ maximal, then $P_\lambda(q) \in \Z_{\geq 0}[q]$.
\end{quote}
The proof gives an explicit decomposition $\sigma_1 = \sum_k n_k(\lambda)(1+q)^k q^{(N-k)/2}$ with $n_k(\lambda) \in \Z_{\geq 0}$, summed over $k \equiv N \pmod 2$, where $N = \binom{n}{2}$. Verified on all $\lambda$ with $V_\lambda^H \neq 0$ at $n \le 7$.
\end{abstract}

\section{Setup and main theorem}

Let $\Hq$ be the Hecke algebra of $\Sym_n$ over $\Z[q, q^{-1}]$, with generators $T_1, \ldots, T_{n-1}$ satisfying
\[ T_i^2 = (q-1) T_i + q, \qquad T_i T_{i+1} T_i = T_{i+1} T_i T_{i+1}, \qquad T_i T_j = T_j T_i\, (|i-j|\geq 2). \]

Working in $\Hq \otimes_{\Z[q]} \Z[v, v^{-1}]$ with $q = v^2$ allows us to use the Kazhdan--Lusztig basis. For each partition $\lambda \vdash n$ we have the cell (Specht) module $\Vlam$; it carries a canonical Kazhdan--Lusztig basis $\{b_w : w \in \mathcal{C}_\lambda\}$ indexed by a left cell $\mathcal{C}_\lambda$ in $\Sym_n$, and equivalently by standard Young tableaux of shape $\lambda$. This basis is bar-invariant and exhibits the \emph{$W$-graph positivity} of Kazhdan--Lusztig (in type A proved elementarily; in general by Elias--Williamson).

Fix the staircase reduced word for the longest element $w_0 \in \Sym_n$:
\[ \underline{w}_0 = (i_1, i_2, \ldots, i_N), \qquad N = \binom{n}{2}, \]
where the sequence is $1;\, 2, 1;\, 3, 2, 1;\, \ldots;\, n-1, n-2, \ldots, 1$. Define the \emph{staircase Bott--Samelson operator}
\[ \Pibs := \prod_{j=1}^{N} (T_{i_j} + 1) \in \Hq. \]

\begin{theorem}[$\sigma_1$-G1]
\label{thm:main}
Let $\lambda \vdash n$ and $\sigma_1(\Vlam) := \trace(\Pibs|_{\Vlam}) \in \Z[q]$. Then there exist non-negative integers $n_k(\lambda) \in \Z_{\geq 0}$, for $k$ ranging over $\{0, 1, \ldots, N\}$ with $k \equiv N \pmod 2$, such that
\begin{equation}
\sigma_1(\Vlam) = \sum_{\substack{0 \leq k \leq N \\ k \equiv N \!\!\!\!\pmod 2}} n_k(\lambda) \cdot (1+q)^k \cdot q^{(N-k)/2}. \label{eq:main}
\end{equation}
Consequently, if $\sigma_1(\Vlam) = q^{a_\lambda}(1+q)^{b_\lambda} P_\lambda(q)$ with $(a_\lambda, b_\lambda)$ maximal, then
\[ P_\lambda(q) = \sum_k n_k(\lambda) \cdot (1+q)^{k - b_\lambda} \cdot q^{(N-k)/2 - a_\lambda} \in \Z_{\geq 0}[q]. \]
\end{theorem}

\section{The key decomposition of $T_i + 1$ in the KL basis}

\begin{lemma}[KL decomposition of $T_i + 1$]
\label{lem:decomp}
In the Kazhdan--Lusztig basis $\{b_w\}$ of $\Vlam$, the action of $T_i + 1$ takes the form
\[ T_i + 1 = (1+q)\, D_i + v\, M_i, \]
where $D_i$ and $M_i$ are matrices with non-negative integer entries, satisfying:
\begin{enumerate}[(i)]
    \item $D_i$ is diagonal: $(D_i)_{ww} = [s_i \cdot w < w]$ (indicator that $s_i$ is a left descent of $w$).
    \item $M_i$ is off-diagonal with entries in $\Z_{\geq 0}$, and $(M_i)_{z,w} \neq 0$ implies $\ell(z) - \ell(w)$ is \emph{odd}.
\end{enumerate}
\end{lemma}

\begin{proof}
Recall the Kazhdan--Lusztig element $C'_{s_i} \in \Hq$ acts on $\Vlam$ (in the KL basis) by
\begin{equation}
\label{eq:kl-action}
 C'_{s_i} \cdot b_w = \begin{cases} (v + v^{-1}) b_w & \text{if } s_i \cdot w < w, \\ b_{s_i \cdot w} + \displaystyle\sum_{\substack{z\,:\, s_i \cdot z < z \\ \mu(z, w) \neq 0}} \mu(z, w) b_z & \text{if } s_i \cdot w > w, \end{cases}
\end{equation}
where $\mu(z, w)$ are the Kazhdan--Lusztig $\mu$-coefficients (non-negative integers, by KL positivity). This is the defining action of $C'_{s_i}$ on a $W$-graph cell module.

The relation $T_i + 1 = v\, C'_{s_i}$ (standard, with $q = v^2$ and the Hecke convention $T_i^2 = (q-1)T_i + q$) gives:
\begin{itemize}
    \item If $s_i \cdot w < w$: $(T_i + 1) b_w = v(v + v^{-1}) b_w = (v^2 + 1) b_w = (1+q) b_w$.
    \item If $s_i \cdot w > w$: $(T_i + 1) b_w = v \cdot b_{s_i w} + v \sum_z \mu(z, w) b_z$.
\end{itemize}
Separating the diagonal and off-diagonal parts yields the decomposition $(T_i + 1) = (1+q) D_i + v M_i$ with the stated properties. The entries of $D_i$ are in $\{0, 1\}$ and those of $M_i$ are non-negative integers by KL positivity.

The parity statement in (ii) is the \emph{parity vanishing} of KL polynomials for $\Sym_n$: $\mu(z, w) \neq 0$ implies $\ell(w) - \ell(z)$ is odd (see \cite[\S 6]{KL}, or \cite{EW} for the general Soergel proof). In the case $z = s_i \cdot w$ with $s_i w > w$, we have $\ell(z) = \ell(w) + 1$, clearly odd. \qedhere
\end{proof}

\section{Expanding $\Pibs$ and computing the trace}

Applying Lemma~\ref{lem:decomp} at each factor:
\begin{align}
\Pibs = \prod_{j=1}^{N} \bigl[(1+q) D_{i_j} + v M_{i_j}\bigr] = \sum_{S \subseteq \{1, \ldots, N\}} (1+q)^{|S|} \cdot v^{N - |S|} \cdot \prod_{j=1}^{N} X_j^{S}, \label{eq:expand}
\end{align}
where the outer product respects the order $j = 1, 2, \ldots, N$ and $X_j^S := D_{i_j}$ if $j \in S$, $X_j^S := M_{i_j}$ if $j \notin S$.

\begin{lemma}[Parity of trace]
\label{lem:parity}
For any $S \subseteq \{1, \ldots, N\}$, $\trace\bigl(\prod_{j=1}^{N} X_j^{S}\bigr) \in \Z_{\geq 0}$, and the trace vanishes unless $|S^c| = N - |S|$ is even.
\end{lemma}

\begin{proof}
Each $D_i$ and $M_i$ is a matrix with non-negative integer entries. Hence $\prod_j X_j^S$ has non-negative integer entries, and its trace is in $\Z_{\geq 0}$.

For the parity claim: the $D$-operators are diagonal and preserve length-parity. The $M$-operators, by Lemma~\ref{lem:decomp}(ii), connect vertices of \emph{opposite} length-parity. The diagonal entry of $\prod_j X_j^S$ at vertex $w$ is a sum over closed paths $w = w_0 \to w_1 \to \cdots \to w_{N-1} \to w_N = w$ where, at step $j$, if $j \in S$ the path stays put ($w_j = w_{j-1}$) and the weight is $(D_{i_j})_{w_{j-1}, w_{j-1}}$; if $j \notin S$ the path may move and the weight is $(M_{i_j})_{w_{j-1}, w_j}$. Each $M$-step flips length-parity of the vertex, each $D$-step preserves it. For the path to close (return to $w$), the total number of parity flips must be even, i.e., $|S^c| = N - |S|$ is even. \qedhere
\end{proof}

\begin{proof}[Proof of Theorem~\ref{thm:main}]
From~\eqref{eq:expand} and Lemma~\ref{lem:parity},
\[
\sigma_1(\Vlam) = \trace(\Pibs|_\Vlam) = \sum_{\substack{S \subseteq \{1,\ldots,N\} \\ |S^c| \text{ even}}} (1+q)^{|S|} v^{N - |S|} N_S,
\]
where $N_S := \trace\bigl(\prod_j X_j^S\bigr) \in \Z_{\geq 0}$. Since $|S^c| = N - |S|$ is even, $v^{N - |S|} = q^{(N-|S|)/2}$ is a non-negative integer power of $q$. Grouping by $k := |S|$:
\[
\sigma_1(\Vlam) = \sum_{\substack{0 \leq k \leq N \\ k \equiv N \!\pmod 2}} (1+q)^k \, q^{(N-k)/2} \cdot \underbrace{\sum_{|S| = k} N_S}_{=: n_k(\lambda) \, \in \, \Z_{\geq 0}}.
\]
This is equation~\eqref{eq:main}.

For the non-negativity of $P_\lambda$: the $(1+q)$-adic valuation of $\sigma_1$ is $b_\lambda = \min\{k : n_k(\lambda) > 0\}$, and the $q$-adic valuation is $a_\lambda = \min_k\{(N-k)/2 : n_k(\lambda) > 0\} = (N - k_{\max})/2$ where $k_{\max} = \max\{k : n_k(\lambda) > 0\}$. (These equalities use that the terms $(1+q)^k q^{(N-k)/2}$ are coprime to each other with respect to both $q$ and $1+q$, and that the $n_k$ are non-negative so no cancellation occurs.)

Dividing both sides of~\eqref{eq:main} by $q^{a_\lambda}(1+q)^{b_\lambda}$:
\[
P_\lambda(q) = \sum_{k} n_k(\lambda)\,(1+q)^{k - b_\lambda}\, q^{(N-k)/2 - a_\lambda}.
\]
For every $k$ with $n_k(\lambda) > 0$: $k \geq b_\lambda$ (so $(1+q)^{k - b_\lambda}$ is a non-negative polynomial) and $(N-k)/2 \geq a_\lambda$ (so $q^{(N-k)/2 - a_\lambda}$ is a non-negative polynomial). Each summand is a non-negative polynomial in $q$, hence so is $P_\lambda$. \qedhere
\end{proof}

\section{Verification and example}

\begin{example}[$\lambda = (3,1)$ at $n = 4$]
Here $N = 6$ and $\sigma_1 = q(1+q)^2(q^2 + 4q + 1) = q + 6q^2 + 10q^3 + 6q^4 + q^5$. Parity $N \equiv 0 \pmod 2$ forces $k$ even. Solving the linear system~\eqref{eq:main} gives
\[
n_0 = 0, \quad n_2 = 2, \quad n_4 = 1, \quad n_6 = 0,
\]
so $\sigma_1 = 2\, q^2 (1+q)^2 + q (1+q)^4$, each term manifestly in $\Z_{\geq 0}[q]$. Consequently $b_\lambda = 2$, $a_\lambda = 1$, and
\[
P_\lambda = 2 q + (1+q)^2 = q^2 + 4q + 1 \in \Z_{\geq 0}[q]. \quad\checkmark
\]
\end{example}

\begin{example}[$\lambda = (3,2,1)$ at $n = 6$]
$N = 15$ forces $k$ odd. Solving gives
\[
n_1 = 2, \quad n_3 = 15, \quad n_5 = 9, \quad n_7 = 1,
\]
all other $n_k = 0$. So $b_\lambda = 1$, $a_\lambda = (15 - 7)/2 = 4$, and
\[
P_\lambda = 2 q^3 + 15 q^2 (1+q)^2 + 9 q (1+q)^4 + (1+q)^6 = q^6 + 15 q^5 + 66 q^4 + 106 q^3 + 66 q^2 + 15 q + 1. \quad\checkmark
\]
\end{example}

\begin{example}[Degree-$10$ novel orphan at $n=7$]
For $\lambda = (3,2,1,1) \vdash 7$ at $k=2$, one of the atoms of the factored $\sigma_2$ core is the irreducible degree-$10$ palindrome $1 + 16q + 105q^2 + 369q^3 + 759 q^4 + 961 q^5 + \cdots$ with $P(1) = 3461$ prime (Apr 24 work). Theorem~\ref{thm:main} gives a non-negativity proof of $P_\lambda$ for $\sigma_1$ only, so this example is \emph{not} covered here; it remains open whether the analog $\sigma_k$-G1 holds. The novel orphan lives at $k = 2$, outside the scope of the present theorem.
\end{example}

\begin{remark}[What is the scope?]
Theorem~\ref{thm:main} covers $\sigma_1$ for \emph{all} $\lambda$, both rank-1 and higher rank. We do \emph{not} handle $\sigma_k$ for $k \geq 2$. The key ingredient that fails for $k \geq 2$: the exterior power $\Lambda^k(\Pibs)$ acting on $\Lambda^k \Vlam$ does not admit the clean $D + vM$ decomposition. An alternative argument (e.g., via suitable Soergel multi-bimodules or tensor-power extensions) is required.
\end{remark}

\begin{remark}[What is used, what is needed]
The proof uses:
\begin{enumerate}[(a)]
    \item The existence of the Kazhdan--Lusztig basis of $\Vlam$ with the $W$-graph action~\eqref{eq:kl-action}. This exists over $\Z[v, v^{-1}]$ in type A.
    \item $W$-graph positivity: the $\mu$-coefficients are non-negative integers. For $\Sym_n$ this is elementary; in general (Elias--Williamson).
    \item Kazhdan--Lusztig parity vanishing: $\mu(z, w) \neq 0 \Rightarrow \ell(w) - \ell(z)$ odd.
\end{enumerate}
No semisimplicity or genericity of $q$ is required. The theorem holds with the same proof over any ring where the Hecke algebra and cell module are defined and the above positivity statements hold.
\end{remark}

\section{Computational verification}

The decomposition~\eqref{eq:main} was verified computationally for all $\lambda \vdash n$ with $V_\lambda^H \neq 0$ at $n \leq 7$:
\begin{itemize}
    \item $n = 3$: $\lambda \in \{(3), (2,1)\}$. All $n_k \geq 0$.
    \item $n = 4$: $\lambda \in \{(4), (3,1), (2,2)\}$. All $n_k \geq 0$.
    \item $n = 5$: $\lambda \in \{(5), (4,1), (3,2), (3,1,1), (2,2,1)\}$. All $n_k \geq 0$.
    \item $n = 6$: $\lambda \in \{(6), (5,1), (4,2), (4,1,1), (3,3), (3,2,1), (2,2,2)\}$. All $n_k \geq 0$.
    \item $n = 7$: $\lambda \in \{(7), (6,1), (5,2), (5,1,1), (4,3), (4,2,1), (4,1,1,1), (3,3,1), (3,2,2), (3,2,1,1), (2,2,2,1)\}$. All $n_k \geq 0$.
\end{itemize}
For $\lambda$ with $V_\lambda^H = 0$ (e.g., $\lambda = (2, 1^{n-2})$ at $n \geq 4$), $\sigma_1 = 0$ trivially and all $n_k = 0$.

Scripts: \texttt{\~{}/projects/scratch/2026-04-24-sigma1-G1/\{compute\_Pi\_integral.py, verify\_proof.py\}}.

\section{Discussion}

\subsection*{What this solves}
This resolves the G1 question for $\sigma_1$ (the trace of $\Pibs$ on each irreducible $V_\lambda$). Prior to today, the known properties of $\sigma_1$ were:
\begin{itemize}
    \item $\sigma_1 \in \Z[q]$ (integrality, from cell-module integral form).
    \item $\sigma_1$ is palindromic of degree $N$ (from bar-involution, Apr 23).
    \item $\sigma_1 \in \Z_{\geq 0}[q]$ (from KL positivity, by taking trace of a non-negative matrix).
    \item $\sigma_1(q) > 0$ for all $q > 0$ (eigenvalue positivity, proved Apr 16).
\end{itemize}
These do \emph{not} force the non-negativity of the core $P_\lambda$ after extracting $q^{a_\lambda}(1+q)^{b_\lambda}$: the palindromic cyclotomic $\Phi_6(q) = q^2 - q + 1$ satisfies all four properties yet has a negative coefficient. So G1 requires genuinely new structure.

The new structure is the decomposition~\eqref{eq:main}: the $(1+q)$-factors and $q$-factors appearing in $\sigma_1$ are \emph{not} artifacts of extraction; they are built in, one factor of $(1+q)$ per ``$D$-step'' (descent position) and one factor of $q$ per ``two $M$-steps'' (KL-graph traversal).

\subsection*{What this does not solve}
This does not address:
\begin{itemize}
    \item $\sigma_k$-G1 for $k \geq 2$ (exterior power cores). The novel degree-10 orphan palindrome at $n=7$, $\lambda = (3, 2, 1, 1)$, $k = 2$, is \emph{not} covered.
    \item The atom alphabet question: are all $P_\lambda$ (for $\sigma_1$) built from a finite set of irreducible palindromic atoms as $n$ grows?
    \item Identification of $b_\lambda$ and $a_\lambda$ with combinatorial invariants of $\lambda$. The empirical rule ``$b_\lambda = \lfloor n/2 \rfloor$ in the rank-$1$ case'' remains conjectural; here it arises as $\min\{k : n_k(\lambda) > 0\}$ which the proof does not compute explicitly.
\end{itemize}

\subsection*{Extension to higher $\sigma_k$}
For $\sigma_k(\Vlam) = \trace(\Lambda^k \Pibs|_{\Lambda^k \Vlam})$, the natural generalization would be: does $\Lambda^k \Pibs$ admit a decomposition analogous to Lemma~\ref{lem:decomp}, with $(1+q)^{\text{descents}}$ and $v^{\text{edges}}$ contributions? The exterior power of the individual factors $(T_i + 1) = (1+q) D_i + v M_i$ is
\[
\Lambda^k(T_i + 1) = \sum_{j=0}^{k} (1+q)^j v^{k-j} \Lambda^j D_i \otimes \Lambda^{k-j} M_i,
\]
but $\Lambda^{k-j} M_i$ need not have non-negative entries in the induced basis (a minor of a non-negative matrix can be negative). This appears to be the precise obstruction to extending the present proof.

\begin{thebibliography}{10}
\bibitem{KL} D.~Kazhdan and G.~Lusztig, \emph{Representations of Coxeter groups and Hecke algebras}, Invent. Math. 53 (1979), 165--184.
\bibitem{EW} B.~Elias and G.~Williamson, \emph{The Hodge theory of Soergel bimodules}, Annals of Math. 180 (2014), 1089--1136.
\bibitem{Mathas} A.~Mathas, \emph{Iwahori--Hecke algebras and Schur algebras of the symmetric group}, Univ. Lecture Series 15, AMS (1999).
\end{thebibliography}

\end{document}
